\chapter{Реляционное время и предел самозапуска проекторного перехода}

Предыдущая проверка временно записал закон быстрого переключения через внешний параметр $t$.
Но ранняя онтология проекта допускает более сильную постановку: время само
является производным состоянием системы. Тогда нельзя объяснять рождение
материи внешним толчком во времени; нужно проверить, способен ли уже
построенный родитель породить одновременно внутренний ход времени и
проекторное рассогласование.

\section{Признанные подходы к внутреннему времени}

Page и Wootters показали, что глобально стационарное состояние может
описывать динамику подсистемы относительно показаний внутренней
подсистемы-часов; см. \path{Phys.Rev.D 27 (1983) 2885}. Rovelli
формализовал динамику как корреляции частичных наблюдаемых; см.
\path{arXiv:gr-qc/0110035}. Connes и Rovelli предложили гипотезу теплового
времени: верное состояние на алгебре наблюдаемых определяет модулярный
однопараметрический поток; см. \path{arXiv:gr-qc/9406019}.

Эти подходы признаны в литературе, но решают разные задачи. Механизм
Page--Wootters требует выделенной подсистемы-часов, запутанности и
глобального ограничения. Модулярный подход строит поток из состояния, но
сам по себе не создаёт необратимость или смену этого состояния.

\section{В проекте уже имеется нетривиальное модулярное состояние}

Том V построил на $M_{18}(\mathbb C)$ состояние
\begin{equation}
 \rho_\beta=\frac{e^{-\beta h_F}}{\Tr e^{-\beta h_F}},
 \qquad \beta>0,
 \label{eq:v6-existing-modular-state}
\end{equation}
и поток
\begin{equation}
 \sigma_s^{\rho_\beta}(A)
 =\rho_\beta^{is}A\rho_\beta^{-is},
 \qquad
 \delta_\rho(A)=i[-\log\rho_\beta,A]
 =i\beta[h_F,A].
 \label{eq:v6-existing-modular-flow}
\end{equation}
Он выделяет положительно-частотный дифференциал цепи и потому является
настоящей конечномерной реализацией части гипотезы теплового времени.

Здесь важно исправить возможное смешение. Нормированный след
$\tau_{300}$ задаёт нормировку действия. Состояние $\rho_\beta$ не является
этим следом и при $\beta>0$ нетривиально. В пределе $\beta=0$ оно становится
$I/18$, и модулярный поток действительно исчезает.

\section{Точный тест проекторной ориентации}

Высота $h_F$ постоянна на каждом трёхмерном семейном блоке:
\begin{equation}
 h_F\big|_{V_i}=c_i I_3.
\end{equation}
Пусть $P=nn^T$ есть любой ранг-один проектор внутри такого блока. Тогда
\begin{equation}
 [h_F,P]=0,
 \qquad
 \sigma_s^{\rho_\beta}(P)=P.
 \label{eq:v6-modular-flow-fixes-projective-axis}
\end{equation}
Следовательно, существующий модулярный поток различает вершины цепи и
ориентирует межвершинные стрелки, но не вращает проекторную ось
$P\in\mathbb{RP}^2$ и не создаёт доменное рассогласование.

\begin{theorem}[Разделение внутреннего времени и проекторного переключения]
Состояние $\rho_\beta$ создаёт нетривиальный модулярный поток на
межвершинных операторах, но действует тождественно на всей проекторной
орбите внутри каждого семейного триплета. Поэтому уже построенное внутреннее
время не является генератором проекторного переключения.
\end{theorem}

Машинный аудит проверяет это на случайных проекторах $nn^T$: максимальный
остаток модулярного движения равен нулю с машинной точностью, тогда как
межвершинная матричная единица движется нетривиально.

\section{Почему состояние не может само себя столкнуть}

Для любого верного состояния выполнено
\begin{equation}
 \sigma_s^\rho(\rho)=\rho,
\end{equation}
поскольку $\rho$ коммутирует со своими степенями. В конечномерном факторе
модулярный поток внутренний, унитарен и обратим. Он сохраняет спектр
$\rho$ и энтропию фон Неймана.

Значит, собственный модулярный поток состояния упорядочивает изменение
наблюдаемых относительно этого состояния, но не создаёт переход
\begin{equation}
 \rho_{\rm sym}\longrightarrow\rho_{\rm broken}.
\end{equation}
Для такого перехода требуется либо семейство неравновесных состояний и
относительный модулярный cocycle, либо вторая подсистема-час и глобальное
ограничение Page--Wootters, либо строго выведенное открытое/диссипативное
описание.

\begin{nogocriterion}[Запрет самопорождающегося модулярного переключения]
Нельзя одновременно определить время собственным модулярным потоком
$\rho$ и приписать тому же потоку изменение самого $\rho$. Это круговое
объяснение. Необратимость и быстрое переключение должны возникнуть из дополнительных
реляционных данных, но не из внешней функции времени, вставленной вручную.
\end{nogocriterion}

\section{Кандидат внутренних часов из багажа проекта}

В Томе V построен четырёхтактный унитарный сдвиг с фазами
\begin{equation}
 1\longmapsto i\longmapsto-1\longmapsto-i\longmapsto1.
\end{equation}
Он является естественным кандидатом конечных внутренних часов. Однако
текущий результат о резонаторе доказывает лишь задержку и рассеяние. Пока
не предъявлены:
\begin{enumerate}
 \item разложение на часы и проекторную систему;
 \item глобальное стационарное ограничение
 $H_{\rm clock}+H_{\rm sys}=0$;
 \item запутанное состояние, условные срезы которого дают эволюцию
 $(\lambda,P)$;
 \item предел, в котором четыре такта достаточно хорошо аппроксимируют
 монотонные часы периода рождения дефекта.
\end{enumerate}
Поэтому четырёхтактная петля пока подобной часам, но ещё не является
Page--Wootters-часами.

\section{Статус совместного рождения времени и материи}

Гипотеза стала уже и строже. Проект действительно содержит признанный
механизм внутреннего модулярного времени, поэтому внешнее $t$ предыдущего
проверки должно считаться временной координатой вычислительной модели, а не
фундаментальным входом. Но существующий поток не действует на
$\mathbb{RP}^2$-ориентацию и не порождает собственного быстрого переключения.

Совместное рождение времени и материи остаётся открытой гипотезой. Следующий
дешёвый тест должен проверить, можно ли использовать уже доказанную
четырёхтактную петлю как реляционные часы для проекторной системы без новых
состояний и получить условную эволюцию, оставляющую пару $+15/-15$.

\begin{protocol}
\begin{enumerate}
 \item признанная литература по реляционному времени: \textbf{есть};
 \item признанная литература по модулярному времени: \textbf{есть};
 \item проектное верное нетривиальное состояние $\rho_\beta$:
 \textbf{уже построено};
 \item нетривиальный поток межвершинных стрелок: \textbf{да};
 \item движение проекторной оси $P\in\mathbb{RP}^2$ этим потоком:
 \textbf{нет};
 \item изменение $\rho_\beta$ его собственным потоком: \textbf{нет};
 \item необратимость и релаксация из конечного модулярного потока:
 \textbf{нет};
 \item внешний параметр времени предыдущей проверки:
 \textbf{только вычислительная координата};
 \item четырёхтактная петля как готовые внутренние часы:
 \textbf{ещё не доказано};
 \item совместное рождение времени и материи: \textbf{не доказано};
 \item следующая проверка:
 \textbf{реляционные часы четырёхтактной петли и проекторное быстрое переключение}.
\end{enumerate}
\end{protocol}

Машинный сертификат сохранён в
\path{s2t_v6_projective_quench_parent_dynamics_gate_results.json}.